\documentclass[aspectratio=169]{beamer}

\usepackage{cmap}
\usepackage{tikz}
\usepackage{array}
\usepackage{multicol}
\usepackage{booktabs}
\usepackage{csquotes}
\usepackage{amssymb,amsfonts,amsmath}
\usepackage{graphicx}

\usecolortheme{beaver}
\usefonttheme[onlymath]{serif}
\setbeamertemplate{navigation symbols}{}
\setbeamertemplate{footline}[page number]
\setbeamertemplate{blocks}[rounded=true,shadow=true]
\setbeamersize{text margin left=0.5cm,text margin right=0.5cm}

\setbeamerfont{author}{size=\normalsize}
\setbeamerfont{institute}{size=\normalsize}
\setbeamerfont{date}{size=\normalsize}


\title{PersonalGen: Knowledge-Preserving Personalization of\\
Text-to-Image Diffusion Models}

\author{
    Karina Peskova\\
    Consultant: Name Surname, PhD/DSc\\
    Advisor: Name Surname, PhD/DSc
}

\institute{
    Diploma Research Project\\
    Lomonosov Moscow State University
}

\date{2026}


\begin{document}


%------------------------------------------------

\begin{frame}
    \thispagestyle{empty}
    \maketitle
\end{frame}


%------------------------------------------------

\begin{frame}{Personalization of Diffusion Models}

Text-to-image diffusion models demonstrate impressive generation quality,
but adapting them to new concepts remains challenging.

\begin{block}{Goal}
Investigate how personalized concepts are represented inside diffusion models
and develop a knowledge-preserving personalization pipeline.
\end{block}


\begin{block}{Problem}
Fine-tuning methods such as DreamBooth LoRA can learn new concepts,
but may overwrite previously acquired knowledge.
\end{block}


\begin{block}{Solution}
We combine personalization and model editing approaches:

\begin{enumerate}[1)]
    \item Train a personalized LoRA adapter.
    \item Analyze cross-attention representations.
    \item Apply AlphaEdit-based constrained editing.
\end{enumerate}

\end{block}

\end{frame}


%------------------------------------------------

\begin{frame}{Research Motivation}

\begin{columns}

\begin{column}{0.55\textwidth}

Modern diffusion models contain billions of parameters and store complex
visual knowledge.

The main research questions:

\begin{enumerate}[1)]
    \item Where is personalized concept information stored?
    \item Are cross-attention layers sufficient for concept transfer?
    \item How can we preserve pretrained knowledge?
\end{enumerate}

\end{column}


\begin{column}{0.45\textwidth}

\begin{figure}
\centering
\includegraphics[width=0.85\textwidth]{example-image}
\caption{Text-to-image diffusion pipeline}
\end{figure}

\end{column}

\end{columns}

\end{frame}


%------------------------------------------------

\begin{frame}{Problem Statement}


Given a pretrained diffusion model:

\[
M_0
\]

we want to obtain a personalized model:

\[
M_p
\]

which generates a new concept while preserving the original model behavior.


The challenge:

\[
M_p = M_0 + \Delta W
\]

where the update $\Delta W$ should introduce new knowledge without
destroying existing information.


\end{frame}


%------------------------------------------------

\begin{frame}{Related Approaches}


\begin{columns}

\begin{column}{0.5\textwidth}

\begin{block}{DreamBooth}
\begin{itemize}
    \item Subject-driven generation
    \item Full model fine-tuning
    \item High personalization quality
\end{itemize}
\end{block}


\begin{block}{LoRA}
\begin{itemize}
    \item Parameter-efficient adaptation
    \item Low-rank updates
    \item Faster training
\end{itemize}
\end{block}

\end{column}


\begin{column}{0.5\textwidth}

\begin{block}{AlphaEdit}
\begin{itemize}
    \item Knowledge-preserving editing
    \item Controlled parameter modification
    \item Reduced interference
\end{itemize}
\end{block}

\end{column}

\end{columns}


\end{frame}


%------------------------------------------------

\begin{frame}{Proposed Pipeline}


\begin{enumerate}[1)]

\item \textbf{LoRA personalization}

Train Stable Diffusion on a small set of images.

\vspace{0.2cm}

\item \textbf{Attention extraction}

Extract cross-attention key/value matrices.

\vspace{0.2cm}

\item \textbf{Representation analysis}

Construct text embedding matrices:

\begin{itemize}
    \item identifier tokens
    \item content words
    \item all valid tokens
\end{itemize}


\item \textbf{AlphaEdit}

Modify parameters with knowledge preservation constraints.


\end{enumerate}


\end{frame}


%------------------------------------------------

\begin{frame}{Experimental Setup}


\begin{block}{Base Model}

Stable Diffusion v1-4

\end{block}


\begin{block}{Personalization}

DreamBooth LoRA

\begin{itemize}
    \item different identifier tokens
    \item different learning rates
    \item different checkpoints
\end{itemize}

\end{block}


\begin{block}{Evaluation}

Comparison of:

\begin{itemize}
    \item original model
    \item LoRA model
    \item AlphaEdit model
\end{itemize}

\end{block}


\end{frame}


%------------------------------------------------

\begin{frame}{Experiments}


\begin{columns}

\begin{column}{0.5\textwidth}

\begin{figure}

\centering
\includegraphics[width=0.8\textwidth]{example-image-a}

\caption{LoRA generations}

\end{figure}


\end{column}


\begin{column}{0.5\textwidth}


\begin{figure}

\centering
\includegraphics[width=0.8\textwidth]{example-image-b}

\caption{Edited model generations}

\end{figure}


\end{column}

\end{columns}


Experiments analyze:

\begin{itemize}
    \item concept similarity;
    \item preservation of pretrained abilities;
    \item influence of editing parameters.
\end{itemize}


\end{frame}


%------------------------------------------------


\begin{frame}{Ablation Studies}


We investigate the influence of:


\begin{enumerate}[1)]

\item Identifier token selection

\item LoRA learning rate

\item Training checkpoint

\item Text representation strategy

\item Interpolation between models


\end{enumerate}


Interpolation:


\[
W(\alpha)=
(1-\alpha)W_{AlphaEdit}
+\alpha W_{LoRA}
\]


\end{frame}


%------------------------------------------------


\begin{frame}{Discussion}


Main observations:


\begin{itemize}

\item Personalized knowledge is distributed across multiple model components.

\item Cross-attention key/value matrices contain important but incomplete
information about the concept.

\item Stronger personalization may increase interference with pretrained
knowledge.

\end{itemize}


\end{frame}


%------------------------------------------------


\begin{frame}{Conclusion}


\begin{enumerate}[1)]

\item We study knowledge-preserving personalization of diffusion models.

\item We combine DreamBooth LoRA with AlphaEdit-based editing.

\item We analyze the role of cross-attention parameters in concept transfer.

\item Experiments reveal limitations of editing only a subset of model
parameters.

\end{enumerate}


\end{frame}


\end{document}